\subsection{3. Mục tiêu đề tài}

\indent\indent Đề tài hướng tới việc xây dựng một hệ thống học sâu gọn nhẹ, hoạt động theo cơ chế đa nhiệm nhằm giải quyết trọn vẹn bài toán phát hiện và phân loại chữ bị làm giả trong ảnh tài liệu:\vspace{0.3cm}

\textbf{Mục tiêu tổng quát:}
Xây dựng và đánh giá mô hình học sâu đa nhiệm \textbf{LightweightDTD} (~4.15 triệu tham số) kết hợp EfficientNet-B0 và BlockDCT $8 \times 8$ nhằm định vị chính xác mặt nạ pixel giả mạo và phân loại 6 đặc trưng pháp y vi mô.

\textbf{Mục tiêu cụ thể:}
\begin{enumerate}[label=\alph*)]
    \item \textbf{Xử lý và chuẩn hóa dữ liệu:} Xây dựng tập dữ liệu chuẩn hóa Minh-TFR từ tập TFR/DTD-lightweight, lọc bỏ các văn bản phủ định sai lệch và xây dựng tập kiểm thử chuẩn 2,000 mẫu phân tầng.
    \item \textbf{Thiết kế kiến trúc hai nhánh Không gian -- Tần số:} Kết hợp backbone EfficientNet-B0 trích xuất đặc trưng thị giác với nhánh biến đổi BlockDCT $8 \times 8$ trích xuất đặc trưng nén JPEG, hòa trộn qua mô-đun chú ý scSE.
    \item \textbf{Cài đặt bộ giải mã LiteMID đa góc nhìn gọn nhẹ:} Xây dựng bộ giải mã lồng nhau LiteMID sử dụng các khối tích chập Depthwise Separable, tổng hợp đặc trưng đa tỉ lệ từ $H/4$ đến $H/32$.
    \item \textbf{Huấn luyện đa nhiệm (Multi-task Learning):} Thiết lập hàm mất mát liên hợp $\mathcal{L}_{\text{total}} = \mathcal{L}_{\text{seg}} + \lambda \mathcal{L}_{\text{cls}}$ (kết hợp BCE + Dice Loss cho phân đoạn và BCE cho phân loại 6 nhãn).
    \item \textbf{Đánh giá toàn diện:} Đo lường hiệu năng qua các chỉ số Pixel Accuracy, Precision, Recall, IoU, Macro Dice và Top-3 Class Hit Rate.
\end{enumerate}
